\documentclass[conference]{IEEEtran}
\usepackage{iftex}
\ifPDFTeX
  \usepackage[T5]{fontenc}
  \usepackage[utf8]{inputenc}
\else
  \usepackage{fontspec}
  \setmainfont{texgyretermes-regular.otf}[
    BoldFont=texgyretermes-bold.otf,
    ItalicFont=texgyretermes-italic.otf,
    BoldItalicFont=texgyretermes-bolditalic.otf]
\fi
\usepackage{amsmath,booktabs,balance,float}
\title{Giải bài toán $N$ quân hậu bằng SAT và Min-Conflicts}
\author{\IEEEauthorblockN{Đặng Minh Hiếu}
\IEEEauthorblockA{Khoa Công nghệ Thông tin\\Trường Đại học Tôn Đức Thắng}}
\begin{document}
\maketitle
\begin{abstract}
Nghiên cứu biểu diễn bài toán $N$ quân hậu bằng logic mệnh đề,
sinh CNF và giải bằng Glucose3. Hai cách chọn biến trong Min-Conflicts
được đánh giá trên năm kích thước. Việc ưu tiên hậu có nhiều xung đột
nhất giảm số bước trung bình ở $N=32$ và $N=64$, trong khi các lần
khởi động lại làm tăng chi phí trung bình ở một số kích thước nhỏ hơn.
\end{abstract}
\begin{IEEEkeywords}
$N$ quân hậu, logic mệnh đề, CNF, Glucose3, Min-Conflicts.
\end{IEEEkeywords}

\section{Bài toán $N$ quân hậu}

\subsection{Phát biểu bài toán}
Trên bàn cờ $N\times N$, cần đặt đúng $N$ hậu sao cho không có hai
quân cùng hàng, cột hoặc đường chéo. Chỉ số hàng $r$ và cột $c$ bắt
đầu từ $0$. Hai ô khác nhau cùng đường chéo nếu chúng có cùng
$r-c$ hoặc cùng $r+c$.

Với $N=4$, thử đặt hậu ở $(0,2)$. Trên hàng $1$, cột $2$ đã có hậu
và cột $1,3$ nằm chéo với $(0,2)$, nên chọn $(1,0)$. Trên hàng $2$,
các cột $0,2$ đã có hậu và ô $(2,1)$ nằm chéo với $(1,0)$, nên
chọn $(2,3)$. Hàng cuối còn cột $1$ và ô $(3,1)$ không nằm chéo
với ba hậu trước. Các hiệu $r-c$ là $(-2,1,-1,2)$, còn các tổng
$r+c$ là $(2,1,5,4)$. Mỗi dãy có bốn giá trị phân biệt, xác nhận
nghiệm ở Hình~\ref{fig:board}.

\begin{figure}[H]
\centering
\setlength{\tabcolsep}{8pt}
\renewcommand{\arraystretch}{1.1}
\begin{tabular}{r|cccc}
$r\backslash c$ & 0 & 1 & 2 & 3\\ \hline
0 & $\cdot$ & $\cdot$ & \textbf{Q} & $\cdot$\\
1 & \textbf{Q} & $\cdot$ & $\cdot$ & $\cdot$\\
2 & $\cdot$ & $\cdot$ & $\cdot$ & \textbf{Q}\\
3 & $\cdot$ & \textbf{Q} & $\cdot$ & $\cdot$
\end{tabular}
\caption{Nghiệm giải tay trên bàn $4\times4$.}
\label{fig:board}
\end{figure}

\subsection{Cài đặt và tối ưu hóa}
\textit{Mô hình logic mệnh đề.}
Biến Boolean $x_{r,c}$ đúng khi ô $(r,c)$ chứa hậu. Mỗi hàng có
mệnh đề $\bigvee_{c=0}^{N-1}x_{r,c}$ để buộc có ít nhất một hậu.
Chương trình nhóm các ô theo $r$, $c$, $r-c$ và $r+c$. Với mỗi cặp
ô $\{a,b\}$ trong cùng nhóm, điều kiện $\neg(x_a\land x_b)$ trở
thành $\neg x_a\lor\neg x_b$ theo luật De Morgan. Hội của các
mệnh đề tạo thành CNF~\cite{lec09}:
\begin{align}
F_N &= \bigwedge_{r=0}^{N-1}
      \left(\bigvee_{c=0}^{N-1}x_{r,c}\right)\notag\\
    &\quad {}\land\bigwedge_{\{a,b\}\in\mathcal{P}}
      (\neg x_a\lor\neg x_b),
\end{align}
với $\mathcal{P}$ là tập các cặp ô tấn công nhau. Mỗi hàng có
đúng một hậu, nên có $N$ hậu. Ràng buộc cột và chéo loại các cặp
tấn công còn lại. Theo chiều ngược lại, mọi bàn cờ hợp lệ thỏa
từng mệnh đề của $F_N$.

Mã DIMACS của ô $(r,c)$ là $rN+c+1$. Với $N=4$, hàng đầu sinh
\texttt{[1,2,3,4]}. Cặp cùng hàng, cùng cột và cùng chéo lần lượt
có các ví dụ \texttt{[-1,-2]}, \texttt{[-1,-5]}, \texttt{[-1,-6]}.
Mỗi họ chéo có các đường dài $1,2,3,4,3,2,1$, tạo $14$ mệnh đề.
Tổng là $4+24+24+14+14=80$. Tổng quát, mỗi hàng và cột có
$\binom{N}{2}$ cặp. Mỗi họ chéo có một đường dài $N$ và hai
đường cho mỗi độ dài $k<N$. Dùng
$\sum_{k=1}^{N-1}\binom{k}{2}=\binom{N}{3}$, số mệnh đề
cấm đôi của mỗi họ chéo là $\binom{N}{2}+2\binom{N}{3}$.
Cộng các nhóm và $N$ mệnh đề dương thu được
\begin{equation}
C(N)=N+N^2(N-1)+\frac{N(N-1)(2N-1)}{3}.
\end{equation}
Gom nhóm cần $O(N^2)$ thời gian. Sinh và lưu tất cả mệnh đề cấm
đôi cần $O(N^3)$.

\textit{Giải bằng Glucose3.}
Lớp \texttt{NQueens} biểu diễn bàn cờ và kiểm tra nghiệm.
\texttt{CNFEncoder} sinh mệnh đề và xuất DIMACS. \texttt{SATSolver}
nạp CNF vào Glucose3 của PySAT, gọi \texttt{solve()} và đọc model
bằng \texttt{get\_model()} nếu khả thỏa. Với literal dương $\ell$,
thương và dư của phép chia $\ell-1$ cho $N$ cho hàng và cột.
Nghiệm giải tay tương ứng các literal dương $3,5,12,14$.
Ví dụ $12-1=2\cdot4+3$ xác định ô $(2,3)$. Glucose3 trả về đúng
nghiệm này ở lần chạy đầu tiên với $N=4$. Chương trình kiểm tra
nghiệm trước khi in hoặc vẽ bàn cờ.

\textit{Min-Conflicts.}
\texttt{MinConflictsSolver} dùng biến $q_c$ lưu hàng của hậu ở
cột $c$, với miền $\{0,\ldots,N-1\}$. Gọi $g_c(t;q)$ là số hậu
ở các cột khác tấn công ô $(t,c)$. Số cặp xung đột là
\begin{equation}
P(q)=\frac12\sum_{c=0}^{N-1} g_c(q_c;q),
\end{equation}
vì mỗi cặp được tính ở cả hai cột. Nghiệm đạt được khi $P(q)=0$.
Theo Lecture~08~\cite{lec08}, bản chuẩn bắt đầu từ một phép
gán đầy đủ rồi chọn đều ngẫu nhiên một cột có $g_c(q_c;q)>0$.
Bản đề xuất chọn cột có $g_c(q_c;q)$ lớn nhất và phá hòa ngẫu
nhiên. Cả hai gỡ hậu, thử mọi hàng rồi chọn ngẫu nhiên một hàng
có $g_c(t;q)$ nhỏ nhất. Chỉ các cặp liên quan đến hậu được dời
có thể thay đổi. Hàng cũ là một ứng viên, nên $P(q)$ không
tăng sau bước sửa nhưng có thể giữ nguyên.

Ba mảng đếm theo hàng và hai họ đường chéo cho phép gỡ hoặc đặt
hậu trong $O(1)$. Điểm của một hậu đang đặt là tổng ba bộ đếm
tại ô đó trừ $3$ để loại chính nó. Sau khi gỡ hậu, chi phí ô ứng
viên bằng tổng trực tiếp. Mỗi bước quét $N$ quân và $N$ hàng, tốn
$O(N)$ thời gian và $O(N)$ bộ nhớ. Khi hết ngân sách một lần chạy,
chương trình có thể khởi tạo lại. Hết toàn bộ ngân sách được ghi
là chưa tìm được nghiệm.

\textit{Mô phỏng bằng tay.}
Từ $q=(1,3,3,2)$, điểm theo cột là $(1,1,3,1)$. Hậu cột $2$
cùng hàng với hậu cột $1$, cùng chéo với hậu cột $0$ và $3$.
Heuristic chọn cột $2$. Sau khi gỡ hậu, chi phí thử các hàng là
$(0,2,2,3)$. Hàng $0$ cho nghiệm $q=(1,3,0,2)$ sau một bước.
Chương trình trả đúng kết quả khi nhận cùng trạng thái đầu.

Bản chuẩn có thể chọn cột $0$ trước. Chi phí $(0,1,2,2)$ đưa
trạng thái về $(0,3,3,2)$. Một đường đi tiếp theo là chọn cột
$2$ với chi phí $(1,1,3,2)$ và phá hòa chọn hàng $0$, thu được
$(0,3,0,2)$. Chọn cột $0$ với chi phí $(1,0,3,1)$ đưa về
nghiệm $(1,3,0,2)$. Nếu bản chuẩn chọn cột $2$ ngay đầu, nó
cũng có thể kết thúc sau một bước.

\subsection{Thí nghiệm và kết quả}
Thí nghiệm chạy trên Windows 11, Python 3.12.6 và \texttt{python-sat}
1.9.dev17. Mỗi cấu hình chạy $20$ lượt cho mỗi
$N\in\{4,8,16,32,64\}$. Mỗi lượt SAT tạo bộ giải mới. Bảng~\ref{tab:sat}
đo riêng thời gian \texttt{solve()} sau khi nạp CNF. Min-Conflicts
dùng cùng tập seed $2026$--$2045$ và cùng chuỗi bàn khởi tạo, kể
cả sau restart. Mỗi lần khởi tạo được tối đa $10\,000$ bước,
với tối đa $4$ lần khởi động lại. Thời gian Min-Conflicts bao
gồm khởi tạo và restart. Một lượt gán lại hàng cũ vẫn tính
là một bước. Tất cả cấu hình đạt $20/20$ nghiệm ở mỗi $N$.

\begin{table}[H]
\centering
\caption{Số mệnh đề và thời gian Glucose3, trung bình 20 lượt}
\label{tab:sat}
\begin{tabular}{rrrr}
\toprule
$N$ & Số biến & Số mệnh đề & \texttt{solve()} (ms)\\
\midrule
4 & 16 & 80 & 0.0062\\
8 & 64 & 736 & 0.0344\\
16 & 256 & 6336 & 1.0878\\
32 & 1024 & 52608 & 12.9204\\
64 & 4096 & 428800 & 152.6832\\
\bottomrule
\end{tabular}
\end{table}

\begin{table}[H]
\centering
\caption{Kết quả trung bình của hai biến thể Min-Conflicts}
\label{tab:local}
\begin{tabular}{rrrrr}
\toprule
& \multicolumn{2}{c}{Thời gian (ms)} & \multicolumn{2}{c}{Số bước}\\
\cmidrule(lr){2-3}\cmidrule(lr){4-5}
$N$ & Chuẩn & Max & Chuẩn & Max\\
\midrule
4 & 0.0416 & 0.0447 & 7.40 & 7.55\\
8 & 3.5220 & 5.8855 & 1103.15 & 1575.45\\
16 & 0.5442 & 6.1647 & 95.60 & 1084.65\\
32 & 0.9424 & 0.8170 & 97.15 & 77.40\\
64 & 2.2594 & 2.0700 & 127.25 & 112.00\\
\bottomrule
\end{tabular}
\end{table}

\balance
\textit{Kiểm chứng.}
Bảng~\ref{tab:sat} cho thấy số mệnh đề đầu vào tăng từ $80$ lên
$428\,800$ khi $N$ tăng từ $4$ lên $64$. Số mệnh đề khớp với công
thức $C(N)$. Kiểm thử vét cạn $2^{16}$ phép gán ở $N=4$ xác nhận
CNF nhận đúng các bàn cờ hợp lệ. Kiểm thử với ngân sách một
bước xác nhận bước chuyển giải tay của heuristic. Cả $300$
nghiệm thực nghiệm đều được kiểm tra theo hàng, cột và đường chéo.

\textit{Phân tích.} Theo Bảng~\ref{tab:local}, heuristic giảm số bước trung bình
khoảng $20.3\%$ ở $N=32$ và $12.0\%$ ở $N=64$. Theo từng cặp
cùng seed, heuristic dùng ít bước hơn ở $13/20$ và $11/20$ lượt.
Ở $N=16$, trung vị giảm từ $84$ còn $61$ bước, nhưng hai lượt
heuristic phải khởi động lại, làm trung bình tăng từ $95.60$
lên $1084.65$. Chênh lệch giữa trung vị và trung bình cho thấy
các lượt chậm ảnh hưởng lớn tới kết quả. Hai bảng dùng các
phạm vi đo khác nhau. Bảng~\ref{tab:sat} mô tả thời gian riêng
của Glucose3 khi tăng $N$, còn Bảng~\ref{tab:local} so sánh hai
quy tắc chọn biến trong cùng một quy trình tìm kiếm. Số đo thời
gian phản ánh cấu hình Windows đã thử.

\begin{thebibliography}{00}
\bibitem{lec09} N. T. An, \emph{Logical Agents},
Lecture 09, Introduction to Artificial Intelligence,
Ton Duc Thang University, 2026, slide 27.
\bibitem{lec08} N. T. An, \emph{Constraint Satisfaction},
Lecture 08, Introduction to Artificial Intelligence,
Ton Duc Thang University, 2026, slides 29--31.
\end{thebibliography}
\end{document}
